% ============================================================
%  Subfile: Skema Ontologi (figur + tabel kelas/relasi/properti)
%  Included by bab-3.tex via \subfile{onthology}
% ============================================================
\documentclass[main]{subfiles}
\begin{document}
\begin{figure}[htbp]
  \centering
  \begin{tikzpicture}[
  node distance=2cm and 2.5cm,
  % Node Styles
  class/.style={circle, draw=blue!80!black, fill=blue!10, thick, 
    minimum size=2.6cm, align=center, font=\bfseries\scriptsize},
  target/.style={circle, draw=blue!80!black, fill=blue!5, dashed, thick, 
    minimum size=2.6cm, align=center, font=\bfseries\scriptsize},
  prop/.style={rectangle, draw=green!60!black, fill=green!10, thick, 
    minimum height=0.5cm, font=\tiny},
  dtype/.style={rectangle, draw=orange!80!black, fill=orange!10, thick, 
    minimum height=0.5cm, font=\tiny},
  % Edge Styles
  edge/.style={->, >={Stealth[length=2.5mm]}, thick, draw=black!70, font=\scriptsize},
  rel/.style={->, >={Stealth[length=2.5mm]}, thick, draw=blue!70!black},
  cross/.style={->, >={Stealth[length=2.5mm]}, thick, dashed, draw=purple!80!black}
]

  % --- 1. HIERARKI KELAS UTAMA (Backbone) ---
  \node[class] (concept) {Concept};
  \node[class, below right=1.8cm and 2cm of concept] (subtopic) {Subtopic};
  \node[class, below right=1.8cm and 2cm of subtopic] (chapter) {Chapter};
  \node[class, below right=1.8cm and 2cm of chapter] (grade) {Grade/\\Subject};

  \node[target, above left=2cm and 3cm of concept] (target) {ConceptTarget\\(Placeholder)};

  % Relasi Hierarki
  \draw[edge] (grade) -- node[below right] {HAS\_CHAPTER} (chapter);
  \draw[edge] (chapter) -- node[below right] {HAS\_SUBTOPIC} (subtopic);
  \draw[edge] (subtopic) -- node[below right] {HAS\_CONCEPT} (concept);

  % --- 2. RELASI PADA CHAPTER ---
  % Loop dipindah ke atas agar tidak menabrak summary
  \draw[edge, draw=orange!90!black] (chapter) edge[loop above, looseness=4.5]
    node[above, align=center, fill=white, inner sep=1pt, font=\tiny] 
      {NEXT\_CHAPTER \\ PRASYARAT \\ MEMPERSIAPKAN} (chapter);

  % Summary dipindah ke kiri bawah
  \node[prop, below left=0.5cm and 1cm of chapter] (summ) {summary};
  \node[dtype, left=0.3cm of summ] (d_str2) {string};
  \draw[edge] (chapter) -- (summ); \draw[edge] (summ) -- (d_str2);


  % --- 3. DATA PROPERTIES PADA CONCEPT ---
  \node[prop, right=1.5cm of concept] (glos) {glossary\_validated};
  \node[dtype, right=0.3cm of glos] (d_bool) {boolean};
  \draw[edge] (concept) -- (glos); \draw[edge] (glos) -- (d_bool);

  \node[prop, above right=0.8cm and 1cm of concept] (desc) {description};
  \node[dtype, right=0.3cm of desc] (d_str1) {string};
  \draw[edge] (concept) -- (desc); \draw[edge] (desc) -- (d_str1);


  % --- 4. RELASI SEMANTIK (WITHIN-BOOK) ---
  % Digambar 3 panah melengkung sebagai representasi, dilabeli di tengah
  \draw[rel] (concept) to[out=160, in=-20] (target);
  \draw[rel] (concept) to[out=140, in=-40] (target);
  \draw[rel] (concept) to[out=120, in=-60] (target);
  
  % Label diletakkan persis di tengah antara Concept dan Target dengan background putih
  \node[font=\tiny, align=center, fill=white, inner sep=3pt, draw=blue!30, rounded corners] 
    at ($ (concept)!0.5!(target) $) {
    \textbf{15 Relasi Semantik (\textit{Within-Book}):} \\
    MENDEFINISIKAN, MENYEBABKAN, MEMUNGKINKAN, \\
    MENGATUR, BAGIAN\_DARI, TERDIRI\_DARI, \\
    BERGANTUNG\_PADA, BERINTERAKSI\_DENGAN, \dots \\
    \textit{(Daftar lengkap merujuk pada Tabel)}
  };


  % --- 5. RELASI LINTAS-BUKU (CROSS-BOOK) ---
  % Loop diletakkan di sebelah kiri agak ke bawah agar terbebas dari garis lain
  \draw[cross] (concept) edge[loop left, out=200, in=250, looseness=6]
    node[below left, fill=white, inner sep=2pt, font=\tiny, align=left, draw=purple!30, rounded corners] {
      \textbf{Relasi Lintas-Buku:} \\
      LINTAS\_BUKU\_SAMA\_DENGAN \\
      LINTAS\_BUKU\_APLIKASI\_DARI \\
      LINTAS\_BUKU\_PRASYARAT\_UNTUK \\
      LINTAS\_BUKU\_MEMPERDALAM \\
      LINTAS\_BUKU\_BERKAITAN\_DENGAN
    } (concept);

\end{tikzpicture}
  \caption{Topologi Ontologi: Relasi Struktural, Pedagogis, dan Semantik.}
  \label{fig:ontologi}
\end{figure}

\subsubsection*{Definisi Kelas}

Ontologi terdiri dari lima kelas sebagaimana dirangkum pada
Tabel~\ref{tab:kelas-ontologi}.

\begin{table}[htbp]
  \centering
  \renewcommand{\arraystretch}{1.3}
  \caption{Definisi Kelas Ontologi}
  \label{tab:kelas-ontologi}
  \begin{tabularx}{\textwidth}{@{}lX@{}}
    \toprule
    \textbf{Kelas} & \textbf{Deskripsi} \\
    \midrule
    \texttt{Grade/Subject}
      & Tingkat kelas dan mata pelajaran pada Kurikulum Merdeka
        (mis.\ Kelas~XII Fisika). Akar hierarki. \\
    \texttt{Chapter}
      & Bab dalam buku teks, bersumber dari Daftar Isi.
        Memiliki ringkasan akademik (\texttt{summary}). \\
    \texttt{Subtopic}
      & Subbab di bawah \texttt{Chapter}.
        Mengelompokkan \texttt{Concept} dalam satu unit tematik. \\
    \texttt{Concept}
      & Konsep ilmiah yang diekstrak dari teks.
        Bersifat \textit{shared} lintas-dokumen melalui
        \texttt{MERGE} berdasarkan nama. \\
    \texttt{ConceptTarget}
      & \textit{Forward reference}: konsep target yang ada sementara
        sebelum diekstrak penuh; dipromosi menjadi \texttt{Concept}
        pada iterasi berikutnya. \\
    \bottomrule
  \end{tabularx}
\end{table}

\subsubsection*{Relasi Objek (Object Properties)}

Hierarki struktural mengikuti rantai:

\begin{center}
  \texttt{Grade/Subject} $\rightarrow$ \texttt{Chapter}
  $\rightarrow$ \texttt{Subtopic} $\rightarrow$ \texttt{Concept}
\end{center}

\noindent dihubungkan oleh relasi \texttt{HAS\_CHAPTER},
\texttt{HAS\_SUBTOPIC}, dan \texttt{HAS\_CONCEPT}. Pada tingkat
\texttt{Chapter}, terdapat tiga relasi pedagogis tambahan:
\texttt{NEXT\_CHAPTER} (urutan baca), \texttt{PRASYARAT}
(ketergantungan pembelajaran), dan \texttt{MEMPERSIAPKAN} (kebalikan
logis dari \texttt{PRASYARAT}).

Vokabulari relasi semantik antarkonsep dibagi menjadi dua kelompok
dengan peran berbeda dalam \textit{pipeline}:

\begin{enumerate}
  \item \textbf{Relasi dalam-buku (\textit{within-book}):} delapan tipe
        yang ditegakkan oleh \textit{prompt} LLM saat ekstraksi per-bab,
        menangkap struktur konseptual lokal dalam satu mata pelajaran
        (Tabel~\ref{tab:relasi-dalam-buku}).
  \item \textbf{Relasi lintas-buku (\texttt{LINTAS\_BUKU\_*}):} lima
        tipe yang dihasilkan pada tahap \textit{completion}
        lintas-disiplin, menangkap interkoneksi semantik antarmata
        pelajaran yang tidak tersurat dalam satu buku
        (Tabel~\ref{tab:relasi-lintas-buku}).
\end{enumerate}

\begin{longtable}{@{}p{0.32\linewidth}p{0.45\linewidth}p{0.16\linewidth}@{}}
  \caption{Vokabulari Relasi Dalam-Buku
    (\texttt{Concept} $\leftrightarrow$ \texttt{Concept})}
  \label{tab:relasi-dalam-buku}\\
  \toprule
  \textbf{Tipe Relasi} & \textbf{Semantik} & \textbf{Sifat} \\
  \midrule
  \endfirsthead
  \toprule
  \textbf{Tipe Relasi} & \textbf{Semantik} & \textbf{Sifat} \\
  \midrule
  \endhead
  \midrule
  \endfoot
  \bottomrule
  \endlastfoot
  \texttt{BAGIAN\_DARI}
    & A merupakan bagian dari B (komposisi/hierarki)
    & asimetris, hierarkis \\
  \texttt{MENYEBABKAN}
    & A menyebabkan B (kausasi)
    & asimetris \\
  \texttt{BERGANTUNG\_PADA}
    & A bergantung pada B (dependensi fungsional)
    & asimetris \\
  \texttt{MENDEFINISIKAN}
    & A mendefinisikan B (definisional)
    & asimetris \\
  \texttt{MEMPENGARUHI}
    & A mempengaruhi B (pengaruh umum)
    & asimetris \\
  \texttt{BERINTERAKSI\_DENGAN}
    & A dan B berinteraksi resiprokal
    & simetris \\
  \texttt{MEMUNGKINKAN}
    & A memungkinkan B terjadi/berfungsi (\textit{enablement})
    & asimetris \\
  \texttt{MENGATUR}
    & A mengatur atau mengontrol perilaku B (regulasi)
    & asimetris \\
\end{longtable}

\begin{longtable}{@{}p{0.38\linewidth}p{0.42\linewidth}p{0.13\linewidth}@{}}
  \caption{Vokabulari Relasi Lintas-Buku (Lintas-Disiplin)}
  \label{tab:relasi-lintas-buku}\\
  \toprule
  \textbf{Tipe Relasi} & \textbf{Semantik} & \textbf{Sifat} \\
  \midrule
  \endfirsthead
  \toprule
  \textbf{Tipe Relasi} & \textbf{Semantik} & \textbf{Sifat} \\
  \midrule
  \endhead
  \midrule
  \endfoot
  \bottomrule
  \endlastfoot
  \texttt{LINTAS\_BUKU\_SAMA\_DENGAN}
    & Konsep di buku berbeda merujuk pada entitas yang sama
    & simetris \\
  \texttt{LINTAS\_BUKU\_APLIKASI\_DARI}
    & Konsep A adalah penerapan konsep B di disiplin lain
    & asimetris \\
  \texttt{LINTAS\_BUKU\_PRASYARAT\_UNTUK}
    & Konsep A merupakan prasyarat pemahaman konsep B di mapel lain
    & asimetris, hierarkis \\
  \texttt{LINTAS\_BUKU\_MEMPERDALAM}
    & Konsep A memperdalam atau memperluas pemahaman konsep B
    & asimetris \\
  \texttt{LINTAS\_BUKU\_BERKAITAN\_DENGAN}
    & Berkaitan secara umum (\textit{fallback})
    & simetris \\
\end{longtable}

\subsubsection*{Properti Data (Data Properties)}

Setiap node dan edge dalam KG menyimpan properti data yang digunakan
untuk ekstraksi, validasi terminologi, dan komputasi
\textit{embedding}. Tabel~\ref{tab:data-properties} merangkum seluruh
properti beserta domain dan fungsinya.

\begin{table}[htbp]
  \centering
  \renewcommand{\arraystretch}{1.3}
  \caption{Properti Data Ontologi}
  \label{tab:data-properties}
  \begin{tabularx}{\textwidth}{@{}lll X@{}}
    \toprule
    \textbf{Properti} & \textbf{Tipe} & \textbf{Domain} & \textbf{Fungsi} \\
    \midrule
    \texttt{name}
      & \texttt{string} & Semua kelas
      & Nama tampilan node; kunci \texttt{MERGE} pada \texttt{Concept}. \\
    \texttt{grade}
      & \texttt{string} & Semua kelas
      & Atribut tingkat kelas, disalin untuk \textit{filtering} cepat. \\
    \texttt{type}
      & \texttt{string} & \texttt{Grade/Subject}
      & Tipe jenjang (mis.\ SMA, SMP). \\
    \texttt{summary}
      & \texttt{string} & \texttt{Chapter}
      & Ringkasan akademik 2--3 kalimat per bab. \\
    \texttt{chapter}
      & \texttt{string} & \texttt{Subtopic}
      & Nama \texttt{Chapter} induk (denormalisasi untuk kueri cepat). \\
    \texttt{description}
      & \texttt{string} & \texttt{Concept}
      & Deskripsi tekstual konsep; sumber utama \textit{embedding}
        untuk tahap KGC. \\
    \texttt{glossary\_validated}
      & \texttt{boolean} & \texttt{Concept}
      & Status validasi nama konsep terhadap glosarium buku teks. \\
    \texttt{materi\_pokok\_ref}
      & \texttt{string} & \texttt{Concept}
      & Referensi ke Capaian Pembelajaran (CP) Kemendikbudristek
        sebagai \textit{curriculum constraint}. \\
    \midrule
    \texttt{relation\_desc}
      & \texttt{string} & \textit{Edge}
      & Rasionalisasi satu kalimat (Bahasa Indonesia) atas eksistensi
        relasi; dihasilkan LLM dan disimpan bersama setiap edge. \\
    \texttt{relation\_type}
      & \texttt{string} & \textit{Edge}
      & Sub-tipe asli dari prediksi LLM untuk keperluan audit dan
        reprodusibilitas. \\
    \bottomrule
  \end{tabularx}
\end{table}
\end{document}
